\documentclass{leveldoc}
\DocUses{foundations,slingshot,stray_field}
\DocFingerprint{2f8b360ead11c8ec}

\begin{document}

\section{What the level asks}\label{sec:asks}

The game's text:
\leveldescription

One ion is launched from $(\val{xZero}, \val{yZero})$ along $+x$ with $T_0 = \val{Tzero}$, in
Newton's mechanics ($c = \infty$), and must enter the detector
$\val{xDetLo} \le x \le \val{xDetHi}$, $\val{yDetLo} \le y \le \val{yDetHi}$ under each of the
level's four disturbances: a plane wave of amplitude $E_0 = \val{Ezero}$ and frequency
$\omega = \val{omega}$ travelling along $+x$, launched at the phases $0^\circ$, $90^\circ$,
$180^\circ$ and $270^\circ$ of its cycle (\physref{2.3}; every shot is flown under every
disturbance, and the level is solved only if every flight arrives). The player places up to four
charges from the listed magnitudes, either sign, at nodes with
$\val{xRegLo} \le x \le \val{xRegHi}$, $\val{yRegLo} \le y \le \val{yRegHi}$. A flight ends when
the ion enters the detector, touches a charge (radius $\val{R}$), or leaves the world
(\xsecref{foundations}{sec:board}).

The reference places $Q = \val{QOne}$ at $(\val{xOne}, \val{yOne})$ and $Q = \val{QTwo}$ at
$(\val{xTwo}, \val{yTwo})$, above the beam line near the launch and below it near the detector.
The document shows that the hum is exactly a random start height and angle plus a small quiver;
that the reference does not cancel the angle as an imaging lens would, but makes the start's
height irrelevant and folds the downward angles back with its lower charge; and it marches the
four flights numerically.

\section{Setup}\label{sec:setup}

The level file, field by field (the appendix reproduces it):
\setuptable

The ion has charge $q = \val{q}$, mass $m = \val{m}$ and speed $v_0 = \val{vZero}$. The hum
accelerates it by up to $qE_0/m = \val{accel}$; its period is $2\pi/\omega = \val{period}$, and the
crossing to the detector, $t \approx \val{tCross}$, lasts $\val{periods}$ periods. The quiver it
causes has the velocity amplitude $qE_0/(m\omega) = \val{quiverV}$ and the displacement amplitude
$qE_0/(m\omega^2) = \val{quiverX}$. A drift at $\val{quiverV}$ across the beam carries the ion
$\val{driftAt}$ off course by the detector, whose half-height is $\val{halfHeight}$.

\begin{figure}[htbp]
\centering
\begin{tikzpicture}
\begin{axis}[axis equal image, width=0.92\linewidth, xmin=-1.6, xmax=31.6, ymin=-1.6, ymax=21.6,
  xtick={0,5,...,30}, ytick={0,5,...,20}, tick label style={font=\small},
  xlabel={$x$}, ylabel={$y$}, label style={font=\small}]
  \draw[black!8] (0,0) grid (30,20);
  \draw[black!45, dashed] (-1,-1) rectangle (\rawval{xWorld},\rawval{yWorld});
  \draw[black!30, dotted] (\rawval{xRegLo},\rawval{yRegLo}) rectangle (\rawval{xRegHi},\rawval{yRegHi});
  \fill[green!50!black, fill opacity=0.18] (\rawval{xDetLo},\rawval{yDetLo}) rectangle (\rawval{xDetHi},\rawval{yDetHi});
  \draw[green!40!black] (\rawval{xDetLo},\rawval{yDetLo}) rectangle (\rawval{xDetHi},\rawval{yDetHi});
  \addplot[thick, blue!70!black] table {patha.dat};
  \addplot[thick, red!70!black] table {pathb.dat};
  \addplot[thick, teal] table {pathc.dat};
  \addplot[thick, orange!80!black] table {pathd.dat};
  \fill[red!75!black] (\rawval{xOne},\rawval{yOne}) circle[radius=\rawval{R}];
  \fill[red!75!black] (\rawval{xTwo},\rawval{yTwo}) circle[radius=\rawval{R}];
  \fill (\rawval{xZero},\rawval{yZero}) circle[radius=0.15];
\end{axis}
\end{tikzpicture}
\caption{The reference charges (red), the placement region (dotted) and the four flights: the
phases $0^\circ$ (blue), $90^\circ$ (red: it drifts down, passes close to the lower charge and is
thrown back up), $180^\circ$ (teal) and $270^\circ$ (orange).}
\label{fig:board}
\end{figure}

\section{Principles and model}\label{sec:model}

\begin{principles}
  \item Newton's equation of motion \xeqref{foundations}{eq:eom} ($c = \infty$) with the charges'
    Coulomb fields \xeqref{foundations}{eq:coulomb} and the disturbance's wave.
  \item A plane wave, $\vect{E} = E_0\,\uvec{e}\cos(\omega(t - \uvec{k}\cdot\vect{x}/c) + \varphi)$,
    $\vect{B} = \uvec{k}\times\vect{E}/c$, $\uvec{e} = \uvec{z}\times\uvec{k}$, an exact solution
    of Maxwell's equations in vacuum (\physref{2.3}, \Jackson{§7.1}).
  \item The board's rules; the classical Runge--Kutta method, cubic Hermite interpolation and the
    error estimate from the step $2h$; bisection (\xsecref{foundations}{sec:board},
    \xsecref{foundations}{sec:numerics}), as in \xsecref{slingshot}{sec:method} and
    \xsecref{stray_field}{sec:march}.
\end{principles}

With $c = \infty$ the wave's $\uvec{k}\cdot\vect{x}/c$ vanishes and so does $\vect{B}$: the hum is a
field uniform over the board, oscillating in time. Here $\uvec{k} = \uvec{x}$, so
$\uvec{e} = \uvec{y}$ and
\begin{equation}\label{eq:hum}
  \vect{E}_\mathrm{hum}(t) = E_0\cos(\omega t + \varphi)\,\uvec{y}.
\end{equation}
The ion starts at $(\val{xZero}, \val{yZero})$ with $\vect{v} = (v_0, 0)$ at $t = 0$, the wave at
phase $\varphi$, and moves by $m\ddot{\vect{x}} = \vect{F}_Q(\vect{x}) + q\vect{E}_\mathrm{hum}(t)$,
$\vect{F}_Q$ the charges' Coulomb force, until the first of the board's events.

\begin{assumption}{The hum is the same everywhere at once}\label{as:uniform}
The game's world has $c = \infty$, where the wave is exactly \eqref{eq:hum}; nothing is neglected
within the model. A real mains field has a wavelength of thousands of kilometres, so over a lab it
is uniform too; that it oscillates during the flight is the level's simplification (its reality
remarks).
\end{assumption}

\begin{assumption}{Fixed charges, a point ion, no radiation}\label{as:fixed}
The charges stay put and the ion is a point; with $c = \infty$ Liénard's power
\xeqref{foundations}{eq:larmor} vanishes. These are the game's model.
\end{assumption}

\section{Derivation}\label{sec:derivation}

\subsection{The empty board}\label{sec:free}

Without charges the motion separates: $x = x_0 + v_0t$ and $\ddot y = (qE_0/m)\cos(\omega t + \varphi)$
with $y(0) = y_0$, $\dot y(0) = 0$, so
\begin{equation}\label{eq:free}
  y(t) = y_0 + \frac{qE_0}{m\omega^2}\big(\cos\varphi - \cos(\omega t + \varphi)\big)
       - \frac{qE_0}{m\omega}\,\sin\varphi\;t:
\end{equation}
a quiver about a line that starts $\delta y_0 = (qE_0/m\omega^2)\cos\varphi$ off the launch and
drifts with $\delta v_y = -(qE_0/m\omega)\sin\varphi$. Where the line meets the detector or the
world's edge (by bisection on \eqref{eq:free}):
\begin{center}\small\val{free}\end{center}
At $0^\circ$ and $180^\circ$ the hum only shifts the ion, which still arrives; at $90^\circ$ and
$270^\circ$ it tilts its course by $\abs{\delta v_y}/v_0 = \val{quiverV}$ and the ion misses.

\subsection{The quiver and the guiding motion}\label{sec:split}

With the charges the same idea holds exactly. Let $\vect{\xi}(t) = -(qE_0/m\omega^2)\cos(\omega t + \varphi)\,\uvec{y}$,
the quiver of a free ion, which satisfies $m\ddot{\vect{\xi}} = q\vect{E}_\mathrm{hum}$, and write
\begin{equation}\label{eq:split}
  \vect{x}(t) = \vect{x}'(t) + \vect{\xi}(t), \qquad m\ddot{\vect{x}}' = \vect{F}_Q\big(\vect{x}' + \vect{\xi}(t)\big).
\end{equation}
The guiding position $\vect{x}'$ feels only the charges, sampled at the quivering position, and
starts at
\begin{equation}\label{eq:effective}
  y'(0) = y_0 + \frac{qE_0}{m\omega^2}\cos\varphi, \qquad \dot y'(0) = -\frac{qE_0}{m\omega}\sin\varphi,
  \qquad x'(0) = x_0, \quad \dot x'(0) = v_0.
\end{equation}
So the hum is exactly a random start height ($\pm\val{quiverX}$) and angle ($\pm\val{quiverV}$) for
the guiding motion, plus the quiver, at most $\val{quiverX}$, wherever the ion is at the end.
Replacing $\vect{F}_Q(\vect{x}' + \vect{\xi})$ by $\vect{F}_Q(\vect{x}')$ makes the guiding motion a
flight in the charges' static field:
\begin{equation}\label{eq:guiding}
  m\ddot{\vect{x}}' \approx \vect{F}_Q(\vect{x}'), \qquad y_\mathrm{e} \approx y'(t^*) + \xi_y(t^*),
\end{equation}
$t^*$ when $x'$ reaches the entry's $x$. Its error is of first order in $\vect{\xi}\cdot\nabla\vect{F}_Q$
and grows quickly near a charge. The reference's four flights ($d_2$ is the closest approach to
the lower charge; \S\ref{sec:march} for the march):
\begin{center}\small\val{ref}\end{center}
At three phases \eqref{eq:guiding} is good to a few hundredths; at $90^\circ$ the ion passes
$\val{closestB}$ from the lower charge, where a quiver of $\val{quiverX}$ changes the force
substantially, and it is off by $\val{guideWorst}$. The flights must be marched in full.

\subsection{What the two charges do}\label{sec:lens}

Without the hum, the charges map a start $(y_0 + \delta y_0, \dot y = \delta v_y)$ to a height on
the line $x = \val{xDetLo}$. Its derivatives at the plain launch, by central differences of
marched flights,
\begin{equation}\label{eq:lens}
  \pderiv{y_\mathrm{e}}{\,\delta y_0} = \val{dydy}, \qquad \pderiv{y_\mathrm{e}}{\,\delta v_y} = \val{dydv},
\end{equation}
against $1$ and $\val{freeDydv}$ on the empty board. The first is nearly zero: the detector sits
at the focus of the charges for rays that start parallel, so the hum's shift of the start
($0^\circ$, $180^\circ$) hardly matters. The second is not: an imaging lens, which brings every
ray from the source to one point, would make it zero; the reference only halves it. What keeps the
tilted flights on target is the lower charge: a flight aimed down passes nearer to it and is
thrown back up, so the arrival height is lowest, $\val{foldY}$, for $\delta v_y = \val{foldV}$ and
rises on both sides (Figure~\ref{fig:fold}): $\val{yAtMinus}$ at $\delta v_y = -\val{quiverV}$
(passing $\val{closestMinus}$ from the charge), $\val{yAtPlus}$ at $+\val{quiverV}$, and
$\val{yAtSteep}$ already at $-0.12$ (passing $\val{closestSteep}$ from it). Over the hum's range of
angles the heights stay within about two cells, and those slightly below the detector's lower
edge are still climbing and enter through its bottom face.

\begin{figure}[htbp]
\centering
\begin{tikzpicture}
\begin{axis}[width=0.7\linewidth, height=5.2cm, xmin=-0.125, xmax=0.125, ymin=8, ymax=14,
  xlabel={start's transverse velocity $\delta v_y$}, ylabel={height at $x = \val{xDetLo}$},
  tick label style={font=\small}, label style={font=\small},
  xtick={-0.1,-0.05,0,0.05,0.1}, scaled x ticks=false,
  xticklabel style={/pgf/number format/fixed}]
  \fill[green!50!black, fill opacity=0.15] (-0.125,\rawval{yDetLo}) rectangle (0.125,\rawval{yDetHi});
  \addplot[thick, blue!70!black] table {curve.dat};
\end{axis}
\end{tikzpicture}
\caption{The charges' static map: where a flight from the launch reaches the line
$x = \val{xDetLo}$, against its start's transverse velocity. Shaded: the detector's span.}
\label{fig:fold}
\end{figure}

\subsection{The march}\label{sec:march}

The full equations, $m\ddot{\vect{x}} = \vect{F}_Q(\vect{x}) + q\vect{E}_\mathrm{hum}(t)$, are marched
by \xeqref{foundations}{eq:rk4} at $h = \val{H}$ for each phase; the events are located on the
cubic Hermite interpolant by bisection and the error estimated against the step $2h$, as in
\xsecref{stray_field}{sec:march}. The field depends on time, so no energy is conserved; the checks
are the step estimate and, on the empty board, the closed form \eqref{eq:free}, which the march
reproduces to $10^{-9}$.

\section{The reference solution by hand}\label{sec:hand}

\begin{procedure}
  \item\label{st:start} \emph{The phases.} For each $\varphi$, the guiding start \eqref{eq:effective},
    the columns $\delta y_0$ and $\delta v_y$ of the table in \S\ref{sec:split}.
  \item\label{st:march} \emph{The march.} \S\ref{sec:march} for each phase. At every phase the
    ion enters the detector through its \val{faceAll} face, at $t = \val{tA}$, $\val{tB}$,
    $\val{tC}$, $\val{tD}$ and $y = \val{yA}$, $\val{yB}$, $\val{yC}$, $\val{yD}$ for $0^\circ$ to
    $270^\circ$ (step estimates $\val{errA}$, $\val{errB}$, $\val{errC}$, $\val{errD}$).
\end{procedure}

\begin{result}[The reference flights]
At all four phases the ion enters the detector, the closest to an edge $\val{margin}$ inside it.
The level is solved.
\end{result}

\section{Comparison with the game}\label{sec:comparison}

\comparisontable

All differences are at most $\val{largest}$, the game's verified accuracy (Foundations §6); our step
estimates are below $10^{-12}$. The tolerances are ten times the larger of the two computations'
estimates, set before comparing; the empty board's values come from the closed form
\eqref{eq:free}, exact but for rounding.

\section{Other solutions}\label{sec:others}

The level lists four phases; the hum's phase at launch could be any. Marched at
$\val{nSweep}$ phases (every $5^\circ$, $h = 0.01$), the reference arrives at
$\val{nSweepOk}$ of them, at heights from $\val{sweepMin}$ to $\val{sweepMax}$
(Figure~\ref{fig:sweep}); phases that miss: \val{missedPhases}. Its margin is small, though: the
closest of the listed phases arrives $\val{margin}$ inside an edge, and the lowest arrivals of the
sweep enter through the bottom face. The sweep's flights are this document's own; the game
verifies only the four listed phases.

\begin{figure}[htbp]
\centering
\begin{tikzpicture}
\begin{axis}[width=0.7\linewidth, height=4.6cm, xmin=0, xmax=360, ymin=8.5, ymax=11.5,
  xlabel={phase at launch $\varphi$ (degrees)}, ylabel={$y_\mathrm{e}$},
  tick label style={font=\small}, label style={font=\small}, xtick={0,90,180,270,360}]
  \fill[green!50!black, fill opacity=0.15] (0,\rawval{yDetLo}) rectangle (360,\rawval{yDetHi});
  \addplot[only marks, mark size=1.2pt, blue!70!black] table {sweep.dat};
\end{axis}
\end{tikzpicture}
\caption{The reference's entry height against the hum's phase at launch, every $5^\circ$. Shaded:
the detector's span.}
\label{fig:sweep}
\end{figure}

The empty board fails at $90^\circ$ and $270^\circ$ (\S\ref{sec:free}). A setup that cancels the
angle's first-order effect, $\partial y_\mathrm{e}/\partial\,\delta v_y = 0$ (an image of the source
on the detector), would tolerate the hum without relying on a fold; this document does not
search the placements for one.

\section{What the game computes}\label{sec:engine}

The game adds each disturbance's wave to the charges' fields; for $c = \infty$ it is \eqref{eq:hum}
exactly (\physref{2.3}). It integrates Newton's equations with its adaptive eighth-order
Runge--Kutta method (\physref{5}, \physref{6}, \physref{7}), each disturbance's flight verified on
its own. The problem is the same as here, and the differences in the comparison are the two
integrations' errors.

\section{Further reading}\label{sec:reading}

Plane waves: \Jackson{§7.1}. Splitting a motion in an oscillating field into a quiver and a
guiding motion is the starting point of Landau and Lifshitz's \emph{Mechanics}, §30 (there for a
rapidly oscillating field, where the guiding motion's force can be averaged; here the field is
slow and the split is used exactly). Flights under disturbances: \levelref{stray_field}.

\levelfileappendix

\end{document}
